\documentclass[11pt]{article}

\usepackage[T1]{fontenc}
\usepackage[utf8]{inputenc}
\usepackage{lmodern}
\usepackage{amsmath,amssymb,amsthm,mathtools}
\usepackage{microtype}
\usepackage{booktabs}
\usepackage{enumitem}
\usepackage{tabularx}
\usepackage{needspace}
\usepackage{array}
\usepackage[margin=1in,top=0.9in,bottom=1in]{geometry}
\usepackage{fancyhdr}
\usepackage{lastpage}
\usepackage{tikz}
\usepackage[hidelinks]{hyperref}

\pagestyle{fancy}
\fancyhf{}
\lhead{\small BIO 2970 --- Selection Practice Exam S1 Key}
\rhead{\small Answer Key}
\cfoot{\small\thepage\ of \pageref*{LastPage}}
\renewcommand{\headrulewidth}{0.4pt}
\setlength{\headheight}{14pt}

% answer text, for KEY documents only
\newcommand{\keyans}[1]{\par\medskip\noindent\textit{Answer.}\ #1\par\medskip}
\newcommand{\work}[1]{\par\vspace{#1}}
\newcommand{\opt}[1]{\makebox[1.9em][l]{(#1)}}
\newcolumntype{Y}{>{\raggedright\arraybackslash}X}
\setlist[enumerate,1]{leftmargin=2em,labelsep=0.6em}
\hyphenpenalty=1000
\relpenalty=10000
\binoppenalty=10000
\begin{document}

\begin{center}
  {\LARGE\bfseries Natural Selection --- Practice Exam S1 \textit{Answer Key}}\\[5pt]
  {\large Level 1 \quad\textbullet\quad Quiz difficulty}\\[6pt]
  BIO 2970 \quad\textbullet\quad Fitness and Selection, Lectures 23 and 24 \quad\textbullet\quad Fall 2026\\[3pt]
  \rule{\textwidth}{0.8pt}\\[3pt]
  {\small Every answer is explained, including why the tempting wrong choices fail. Grade the answer sheet with the grid below, then read the explanations for anything you missed.}
\end{center}
\vspace{0.4em}

\begin{center}\small\setlength{\tabcolsep}{6.2pt}\begin{tabular}{lcccccccccccc}
\toprule
Question & 1 & 2 & 3 & 4 & 5 & 6 & 7 & 8 & 9 & 10 & 11 & 12 \\
\midrule
Answer & \textbf{C} & \textbf{A} & \textbf{D} & \textbf{B} & \textbf{B} & \textbf{D} & \textbf{A} & \textbf{C} & \textbf{D} & \textbf{C} & \textbf{B} & \textbf{A} \\
\bottomrule\end{tabular}\end{center}
\begin{center}\small\setlength{\tabcolsep}{6.2pt}\begin{tabular}{lcccccccccccc}
\toprule
Question & 13 & 14 & 15 & 16 & 17 & 18 & 19 & 20 & 21 & 22 & 23 & 24 \\
\midrule
Answer & \textbf{C} & \textbf{A} & \textbf{D} & \textbf{B} & \textbf{B} & \textbf{D} & \textbf{A} & \textbf{C} & \textbf{D} & \textbf{C} & \textbf{B} & \textbf{A} \\
\bottomrule\end{tabular}\end{center}
\begin{center}\small\begin{tabular}{lll}\toprule Topic & Questions & Your misses\\\midrule
Fitness, \(w\), \(W\) and \(s\) & 1, 2 & \\
General equation and mean fitness & 13, 21, 23 & \\
Counts or Hardy-Weinberg & 3, 4 & \\
Average excess & 14, 15 & \\
Types of selection & 5, 6, 7, 8 & \\
Special models, \(p\) and \(q\) & 9, 10, 11, 12, 22 & \\
Heterozygote advantage and \(\hat{q}\) & 16, 17 & \\
Heterozygote disadvantage and stability & 24 & \\
Human cases and Darwin & 18, 19 & \\
Drift (Deme lab) & 20 & \\
\bottomrule\end{tabular}\end{center}
\vspace{0.6em}
\begin{enumerate}[label=\textbf{\arabic*.}, leftmargin=2em, itemsep=1.4em]
\needspace{14\baselineskip}
\item Viability is 0.4, 0.4 and 0.2 for AA, Aa and aa. Lifetime fecundity is 5, 5 and 8. What is the absolute fitness \(w\) of aa?
\nopagebreak\par\noindent{\footnotesize\itshape Topic: Fitness, \(w\), \(W\) and \(s\)}
\par\smallskip\noindent\begin{tabularx}{\linewidth}{@{}YYYY@{}}
\makebox[1.9em][l]{(A)}8.2 & \makebox[1.9em][l]{(B)}0.2 & \makebox[1.9em][l]{\textbf{(C)}}\textbf{1.6}\ \(\checkmark\) & \makebox[1.9em][l]{(D)}2.0 \\
\end{tabularx}
\keyans{Absolute fitness is viability times fecundity, offspring left per zygote born. For aa that is 0.2 \(\times\) 8 = 1.6. Adding (8.2) is never right, 0.2 is only the survival part, and 2.0 is the fitness of AA and Aa.}
\needspace{14\baselineskip}
\item Same genotypes. What are the relative fitness \(W\) and the selection coefficient \(s\) of aa?
\nopagebreak\par\noindent{\footnotesize\itshape Topic: Fitness, \(w\), \(W\) and \(s\)}
\par\smallskip\noindent\begin{tabularx}{\linewidth}{@{}YY@{}}
\makebox[1.9em][l]{\textbf{(A)}}\textbf{\(W\) = 0.8, \(s\) = 0.2}\ \(\checkmark\) & \makebox[1.9em][l]{(B)}\(W\) = 0.2, \(s\) = 0.8 \\[3pt]
\makebox[1.9em][l]{(C)}\(W\) = 1.6, \(s\) = \(-\)0.6 & \makebox[1.9em][l]{(D)}\(W\) = 0.8, \(s\) = 0.8 \\
\end{tabularx}
\keyans{AA and Aa have \(w\) = 0.4 \(\times\) 5 = 2.0, the largest. \(W\) = \(w\)/\(w_{\max}\) = 1.6/2.0 = 0.8, and \(s\) = 1 \(-\) \(W\) = 0.2. Relative fitness is never above 1, and \(s\) is how far below the best genotype a genotype sits.}
\needspace{14\baselineskip}
\item A population has 110 AA, 50 Aa and 40 aa. Calculate \(q^2\), the way last year's quiz asked.
\nopagebreak\par\noindent{\footnotesize\itshape Topic: Counts or Hardy-Weinberg}
\par\smallskip\noindent\begin{tabularx}{\linewidth}{@{}YYYY@{}}
\makebox[1.9em][l]{(A)}0.200 & \makebox[1.9em][l]{(B)}0.325 & \makebox[1.9em][l]{(C)}0.447 & \makebox[1.9em][l]{\textbf{(D)}}\textbf{0.106}\ \(\checkmark\) \\
\end{tabularx}
\keyans{Count alleles first. a copies are 2 \(\times\) 40 + 50 = 130 out of 400, so \(q\) = 0.325, and \(q^2\) = 0.106. The 0.200 is the observed aa frequency, 40/200, which is not what ``calculate \(q^2\)'' asks for. 0.325 is \(q\) itself, and 0.447 is the square root of 0.200.}
\needspace{14\baselineskip}
\item Same population. Relative fitnesses are 1, 0.8 and 0.5. Calculate \(\overline{W}\) using \(p^2\), \(2pq\) and \(q^2\).
\nopagebreak\par\noindent{\footnotesize\itshape Topic: Counts or Hardy-Weinberg}
\par\smallskip\noindent\begin{tabularx}{\linewidth}{@{}YYYY@{}}
\makebox[1.9em][l]{(A)}0.85 & \makebox[1.9em][l]{\textbf{(B)}}\textbf{0.86}\ \(\checkmark\) & \makebox[1.9em][l]{(C)}0.77 & \makebox[1.9em][l]{(D)}1.00 \\
\end{tabularx}
\keyans{\(p\) = 0.675, so \(p^2\) = 0.456, \(2pq\) = 0.439 and \(q^2\) = 0.106. \(\overline{W}\) = 0.456(1) + 0.439(0.8) + 0.106(0.5) = 0.859. The 0.85 uses the observed frequencies 0.55, 0.25 and 0.20 instead. The 0.77 is the plain average of 1, 0.8 and 0.5, which ignores how common each genotype is.}
\needspace{14\baselineskip}
\item Lizards of medium body size survive best. Small ones lose fights and large ones are easier for hawks to see. Which type of selection is this?
\nopagebreak\par\noindent{\footnotesize\itshape Topic: Types of selection}
\par\smallskip\noindent\begin{tabularx}{\linewidth}{@{}YYYY@{}}
\makebox[1.9em][l]{(A)}Directional & \makebox[1.9em][l]{\textbf{(B)}}\textbf{Stabilizing}\ \(\checkmark\) & \makebox[1.9em][l]{(C)}Disruptive & \makebox[1.9em][l]{(D)}Balancing \\
\end{tabularx}
\keyans{The middle of the trait is favored and both tails are cut. That is stabilizing selection, a trait word. Balancing is a gene word, about keeping two alleles.}
\needspace{14\baselineskip}
\item In polluted forests the dark peppered moth allele rose from rare to over 90 percent. At the level of the allele, which type of selection is this?
\nopagebreak\par\noindent{\footnotesize\itshape Topic: Types of selection}
\par\smallskip\noindent\begin{tabularx}{\linewidth}{@{}YY@{}}
\makebox[1.9em][l]{(A)}Negative (purifying) & \makebox[1.9em][l]{(B)}Balancing \\[3pt]
\makebox[1.9em][l]{(C)}Stabilizing & \makebox[1.9em][l]{\textbf{(D)}}\textbf{Positive}\ \(\checkmark\) \\
\end{tabularx}
\keyans{A favored allele rising toward fixation is positive selection, a selective sweep. On the color trait itself the same event is directional selection.}
\needspace{14\baselineskip}
\item Finches with very small beaks eat soft seeds and finches with very large beaks crack hard seeds. Birds in between starve. Which type of selection is this?
\nopagebreak\par\noindent{\footnotesize\itshape Topic: Types of selection}
\par\smallskip\noindent\begin{tabularx}{\linewidth}{@{}YYYY@{}}
\makebox[1.9em][l]{\textbf{(A)}}\textbf{Disruptive}\ \(\checkmark\) & \makebox[1.9em][l]{(B)}Stabilizing & \makebox[1.9em][l]{(C)}Directional & \makebox[1.9em][l]{(D)}Balancing \\
\end{tabularx}
\keyans{Both extremes beat the middle, so intermediates are removed and the distribution can split in two.}
\needspace{14\baselineskip}
\item A new frameshift mutation in an essential gene keeps appearing and keeps being removed because carriers rarely reproduce. Which type of selection removes it?
\nopagebreak\par\noindent{\footnotesize\itshape Topic: Types of selection}
\par\smallskip\noindent\begin{tabularx}{\linewidth}{@{}YY@{}}
\makebox[1.9em][l]{(A)}Positive & \makebox[1.9em][l]{(B)}Balancing \\[3pt]
\makebox[1.9em][l]{\textbf{(C)}}\textbf{Negative (purifying)}\ \(\checkmark\) & \makebox[1.9em][l]{(D)}Disruptive \\
\end{tabularx}
\keyans{Removing harmful new mutations is negative, or purifying, selection. It keeps genes from slowly degrading by mutation and drift.}
\needspace{14\baselineskip}
\item In cattle, hornless (H) is completely dominant to horned (h). Horned animals are injured more often and survive less. Which selection model is best to use?
\nopagebreak\par\noindent{\footnotesize\itshape Topic: Special models, \(p\) and \(q\)}
\par\smallskip
\begin{enumerate}[label={}, leftmargin=2.4em, labelwidth=0pt, itemsep=2pt, topsep=0pt]
\item \hspace*{-2.4em}\makebox[2.4em][l]{(A)}Complete dominance with selection against the dominant
\item \hspace*{-2.4em}\makebox[2.4em][l]{(B)}Heterozygote advantage
\item \hspace*{-2.4em}\makebox[2.4em][l]{(C)}No special model applies
\item \hspace*{-2.4em}\makebox[2.4em][l]{\textbf{(D)}}\textbf{Complete dominance with selection against the recessive}\ \(\checkmark\)
\end{enumerate}
\keyans{Only hh animals have horns, so only hh is hurt. The allele selected against is recessive. This is the same setup as last year's yak question.}
\needspace{14\baselineskip}
\item Same cattle. Assign the symbols for the selection model equation.
\nopagebreak\par\noindent{\footnotesize\itshape Topic: Special models, \(p\) and \(q\)}
\par\smallskip\noindent\begin{tabularx}{\linewidth}{@{}YY@{}}
\makebox[1.9em][l]{(A)}H = \(q\) and h = \(p\) & \makebox[1.9em][l]{(B)}Either way gives the same \(\Delta q\) \\[3pt]
\makebox[1.9em][l]{\textbf{(C)}}\textbf{H = \(p\) and h = \(q\)}\ \(\checkmark\) & \makebox[1.9em][l]{(D)}H = \(p\) and h = \(p\) \\
\end{tabularx}
\keyans{\(q\) is always the allele selected against, here h. The special formula has a minus sign in front because \(q\) is the allele being removed, so swapping them gives the wrong answer.}
\needspace{14\baselineskip}
\item Same cattle. The selection coefficient for horned animals is 0.3. What is their relative fitness?
\nopagebreak\par\noindent{\footnotesize\itshape Topic: Special models, \(p\) and \(q\)}
\par\smallskip\noindent\begin{tabularx}{\linewidth}{@{}YYYY@{}}
\makebox[1.9em][l]{(A)}0.3 & \makebox[1.9em][l]{\textbf{(B)}}\textbf{0.7}\ \(\checkmark\) & \makebox[1.9em][l]{(C)}1.3 & \makebox[1.9em][l]{(D)}0.0 \\
\end{tabularx}
\keyans{\(W\) = 1 \(-\) \(s\) = 1 \(-\) 0.3 = 0.7.}
\needspace{14\baselineskip}
\item Same cattle. Is it more appropriate to calculate \(\Delta q\) or \(\hat{q}\)?
\nopagebreak\par\noindent{\footnotesize\itshape Topic: Special models, \(p\) and \(q\)}
\par\smallskip\noindent\begin{tabularx}{\linewidth}{@{}YY@{}}
\makebox[1.9em][l]{\textbf{(A)}}\textbf{\(\Delta q\)}\ \(\checkmark\) & \makebox[1.9em][l]{(B)}\(\hat{q}\) \\[3pt]
\makebox[1.9em][l]{(C)}Both are equally appropriate & \makebox[1.9em][l]{(D)}Neither, because h is recessive \\
\end{tabularx}
\keyans{Selection against the recessive has no stable middle frequency. h simply keeps falling, more and more slowly. \(\hat{q}\) belongs to heterozygote advantage.}
\needspace{14\baselineskip}
\item Three genotypes have fecundities 2, 3 and 4 for AA, Aa and aa, frequencies 0.5, 0.4 and 0.1, and equal survival. What is the frequency of a after selection?
\nopagebreak\par\noindent{\footnotesize\itshape Topic: General equation and mean fitness}
\par\smallskip\noindent\begin{tabularx}{\linewidth}{@{}YYYY@{}}
\makebox[1.9em][l]{(A)}0.300 & \makebox[1.9em][l]{(B)}0.615 & \makebox[1.9em][l]{\textbf{(C)}}\textbf{0.385}\ \(\checkmark\) & \makebox[1.9em][l]{(D)}0.154 \\
\end{tabularx}
\keyans{\(\overline{W}\) = 0.5(2) + 0.4(3) + 0.1(4) = 2.6. \(q'\) = [\(\tfrac{1}{2}\)(0.4)(3) + 0.1(4)] / 2.6 = (0.6 + 0.4)/2.6 = 0.385. The 0.300 is \(q\) before selection. The 0.615 forgets the \(\tfrac{1}{2}\) on the heterozygotes. The 0.154 counts only aa.}
\needspace{14\baselineskip}
\item In \(\Delta q\) = (\(q\)/\(\overline{W}\))\(\cdot \alpha_{a}\), which part determines whether the allele increases or decreases?
\nopagebreak\par\noindent{\footnotesize\itshape Topic: Average excess}
\par\smallskip
\begin{enumerate}[label={}, leftmargin=2.4em, labelwidth=0pt, itemsep=2pt, topsep=0pt]
\item \hspace*{-2.4em}\makebox[2.4em][l]{\textbf{(A)}}\textbf{The average excess \(\alpha_{a}\)}\ \(\checkmark\)
\item \hspace*{-2.4em}\makebox[2.4em][l]{(B)}Mean fitness \(\overline{W}\)
\item \hspace*{-2.4em}\makebox[2.4em][l]{(C)}The allele frequency \(q\)
\item \hspace*{-2.4em}\makebox[2.4em][l]{(D)}The Hardy-Weinberg genotype frequencies
\end{enumerate}
\keyans{\(q\) and \(\overline{W}\) are always positive, so the sign of the average excess alone decides the direction. This was a question on last year's quiz.}
\needspace{14\baselineskip}
\item Hardy-Weinberg, \(q\) = 0.3, relative fitnesses 1, 0.8 and 0.6 for AA, Aa and aa. What is the average excess of a?
\nopagebreak\par\noindent{\footnotesize\itshape Topic: Average excess}
\par\smallskip\noindent\begin{tabularx}{\linewidth}{@{}YYYY@{}}
\makebox[1.9em][l]{(A)}+0.14 & \makebox[1.9em][l]{(B)}\(-\)0.32 & \makebox[1.9em][l]{(C)}0.74 & \makebox[1.9em][l]{\textbf{(D)}}\textbf{\(-\)0.14}\ \(\checkmark\) \\
\end{tabularx}
\keyans{\(\overline{W}\) = 0.49(1) + 0.42(0.8) + 0.09(0.6) = 0.88. An a copy sits with a (probability 0.3, fitness 0.6) or with A (probability 0.7, fitness 0.8), so it experiences 0.18 + 0.56 = 0.74. \(\alpha_{a}\) = 0.74 \(-\) 0.88 = \(-\)0.14. The 0.74 forgets to subtract \(\overline{W}\).}
\needspace{14\baselineskip}
\item Wolves have absolute fitnesses 3, 5 and 4 for B1B1, B1B2 and B2B2. What is the equilibrium frequency of B2?
\nopagebreak\par\noindent{\footnotesize\itshape Topic: Heterozygote advantage and \(\hat{q}\)}
\par\smallskip\noindent\begin{tabularx}{\linewidth}{@{}YYYY@{}}
\makebox[1.9em][l]{(A)}0.333 & \makebox[1.9em][l]{\textbf{(B)}}\textbf{0.667}\ \(\checkmark\) & \makebox[1.9em][l]{(C)}0.500 & \makebox[1.9em][l]{(D)}0.800 \\
\end{tabularx}
\keyans{Divide by 5. \(W\) = 0.6, 1, 0.8, so \(s_1\) = 0.4 for B1B1 and \(s_2\) = 0.2 for B2B2. With \(q\) = freq(B2), \(\hat{q}\) = \(s_1\)/(\(s_1\) + \(s_2\)) = 0.4/0.6 = 0.667. The 0.333 swaps \(s_1\) and \(s_2\). Logic check, B1B1 is hurt more, so B1 should be rarer and B2 more common.}
\needspace{14\baselineskip}
\item Why is the equilibrium under heterozygote advantage stable?
\nopagebreak\par\noindent{\footnotesize\itshape Topic: Heterozygote advantage and \(\hat{q}\)}
\par\smallskip
\begin{enumerate}[label={}, leftmargin=2.4em, labelwidth=0pt, itemsep=2pt, topsep=0pt]
\item \hspace*{-2.4em}\makebox[2.4em][l]{(A)}The heterozygotes cannot reproduce
\item \hspace*{-2.4em}\makebox[2.4em][l]{\textbf{(B)}}\textbf{Whichever allele is rare sits mostly in fit heterozygotes, so it is favored and rises back}\ \(\checkmark\)
\item \hspace*{-2.4em}\makebox[2.4em][l]{(C)}Mutation keeps restoring both alleles
\item \hspace*{-2.4em}\makebox[2.4em][l]{(D)}Drift pushes frequencies toward 0.5
\end{enumerate}
\keyans{A rare allele almost always pairs with the other allele, so the fitness it experiences is the heterozygote's, the highest. Its average excess is positive and it increases. A common allele often lands in its weak homozygote and falls. Both pushes point toward \(\hat{q}\).}
\needspace{14\baselineskip}
\item In a region with intense malaria, which genotype has the highest fitness?
\nopagebreak\par\noindent{\footnotesize\itshape Topic: Human cases and Darwin}
\par\smallskip\noindent\begin{tabularx}{\linewidth}{@{}YY@{}}
\makebox[1.9em][l]{(A)}HbA/HbA & \makebox[1.9em][l]{(B)}HbS/HbS \\[3pt]
\makebox[1.9em][l]{(C)}All three are equal & \makebox[1.9em][l]{\textbf{(D)}}\textbf{HbA/HbS}\ \(\checkmark\) \\
\end{tabularx}
\keyans{Carriers resist malaria and do not get sickle cell disease. HbA/HbA is susceptible to malaria and HbS/HbS has sickle cell disease.}
\needspace{14\baselineskip}
\item Which three things must be true for selection to change a trait across generations?
\nopagebreak\par\noindent{\footnotesize\itshape Topic: Human cases and Darwin}
\par\smallskip
\begin{enumerate}[label={}, leftmargin=2.4em, labelwidth=0pt, itemsep=2pt, topsep=0pt]
\item \hspace*{-2.4em}\makebox[2.4em][l]{\textbf{(A)}}\textbf{It varies, it is heritable, and it affects fitness}\ \(\checkmark\)
\item \hspace*{-2.4em}\makebox[2.4em][l]{(B)}It is dominant, common and visible
\item \hspace*{-2.4em}\makebox[2.4em][l]{(C)}It is new, beneficial and fixed
\item \hspace*{-2.4em}\makebox[2.4em][l]{(D)}It varies, it is recessive, and it is rare
\end{enumerate}
\keyans{These are the three requirements in Session 24. If any one is missing, selection cannot change the next generation.}
\needspace{14\baselineskip}
\item In the Deme lab, many populations start at 0.5 and run for 50 generations. Histogram X is tightly bunched near zero change and histogram Y is spread widely. Which is true?
\nopagebreak\par\noindent{\footnotesize\itshape Topic: Drift (Deme lab)}
\par\smallskip
\begin{enumerate}[label={}, leftmargin=2.4em, labelwidth=0pt, itemsep=2pt, topsep=0pt]
\item \hspace*{-2.4em}\makebox[2.4em][l]{(A)}X came from smaller populations than Y
\item \hspace*{-2.4em}\makebox[2.4em][l]{(B)}X experienced selection and Y did not
\item \hspace*{-2.4em}\makebox[2.4em][l]{\textbf{(C)}}\textbf{X came from larger populations than Y}\ \(\checkmark\)
\item \hspace*{-2.4em}\makebox[2.4em][l]{(D)}Y ran for fewer generations
\end{enumerate}
\keyans{Drift is random sampling of gametes. Smaller populations sample fewer gametes, so chance moves their frequencies further, giving a wider spread.}
\needspace{14\baselineskip}
\item Why is the heterozygote term multiplied by \(\tfrac{1}{2}\) in general equation \#1?
\nopagebreak\par\noindent{\footnotesize\itshape Topic: General equation and mean fitness}
\par\smallskip
\begin{enumerate}[label={}, leftmargin=2.4em, labelwidth=0pt, itemsep=2pt, topsep=0pt]
\item \hspace*{-2.4em}\makebox[2.4em][l]{(A)}Heterozygotes have half the fitness
\item \hspace*{-2.4em}\makebox[2.4em][l]{(B)}Heterozygotes are half as common
\item \hspace*{-2.4em}\makebox[2.4em][l]{(C)}It corrects for Hardy-Weinberg
\item \hspace*{-2.4em}\makebox[2.4em][l]{\textbf{(D)}}\textbf{Only one of a heterozygote's two allele copies is the allele being tracked}\ \(\checkmark\)
\end{enumerate}
\keyans{The numerator counts copies of one allele. Each homozygote contributes both copies and each heterozygote contributes one of its two.}
\needspace{14\baselineskip}
\item Dominant allele B makes beetles conspicuous. BB and Bb have relative fitness 0.6, and bb has 1. The frequency of B is 0.2, Hardy-Weinberg. What is \(\Delta q\) for B?
\nopagebreak\par\noindent{\footnotesize\itshape Topic: Special models, \(p\) and \(q\)}
\par\smallskip\noindent\begin{tabularx}{\linewidth}{@{}YYYY@{}}
\makebox[1.9em][l]{(A)}\(-\)0.013 & \makebox[1.9em][l]{(B)}\(-\)0.051 & \makebox[1.9em][l]{\textbf{(C)}}\textbf{\(-\)0.060}\ \(\checkmark\) & \makebox[1.9em][l]{(D)}\(-\)0.021 \\
\end{tabularx}
\keyans{B is selected against and dominant, so it is \(q\) = 0.2, \(p\) = 0.8, \(s\) = 0.4. Top, \(-\)sp\(^2 q\) = \(-\)0.4 \(\times\) 0.64 \(\times\) 0.2 = \(-\)0.0512. Bottom, 1 \(-\) 2sq + sq\(^2\) = 1 \(-\) 0.16 + 0.016 = 0.856. \(\Delta q\) = \(-\)0.060. The \(-\)0.013 uses the recessive formula, \(-\)0.051 forgets the bottom, and \(-\)0.021 swaps \(p\) and \(q\).}
\needspace{14\baselineskip}
\item Every genotype's fitness is multiplied by 10. What happens to \(\Delta q\)?
\nopagebreak\par\noindent{\footnotesize\itshape Topic: General equation and mean fitness}
\par\smallskip
\begin{enumerate}[label={}, leftmargin=2.4em, labelwidth=0pt, itemsep=2pt, topsep=0pt]
\item \hspace*{-2.4em}\makebox[2.4em][l]{(A)}It becomes 10 times larger
\item \hspace*{-2.4em}\makebox[2.4em][l]{\textbf{(B)}}\textbf{Nothing, because \(\overline{W}\) is multiplied by 10 too and the factor cancels}\ \(\checkmark\)
\item \hspace*{-2.4em}\makebox[2.4em][l]{(C)}It becomes 10 times smaller
\item \hspace*{-2.4em}\makebox[2.4em][l]{(D)}It flips sign
\end{enumerate}
\keyans{Selection only responds to how each fitness compares with the mean. Rescaling every fitness rescales \(\overline{W}\) by the same factor.}
\needspace{14\baselineskip}
\item Heterozygotes have the lowest fitness and both homozygotes are fitter. What happens over many generations?
\nopagebreak\par\noindent{\footnotesize\itshape Topic: Heterozygote disadvantage and stability}
\par\smallskip
\begin{enumerate}[label={}, leftmargin=2.4em, labelwidth=0pt, itemsep=2pt, topsep=0pt]
\item \hspace*{-2.4em}\makebox[2.4em][l]{\textbf{(A)}}\textbf{An unstable equilibrium, so one allele goes to fixation}\ \(\checkmark\)
\item \hspace*{-2.4em}\makebox[2.4em][l]{(B)}Both alleles are kept at \(\hat{q}\)
\item \hspace*{-2.4em}\makebox[2.4em][l]{(C)}Both alleles are lost
\item \hspace*{-2.4em}\makebox[2.4em][l]{(D)}Nothing changes
\end{enumerate}
\keyans{The equilibrium point exists but frequencies run away from it. Translocation heterozygotes, which make unbalanced gametes, are the course's example (Problem Set 12 question 9).}
\end{enumerate}

\newpage
\begin{center}
  {\Large\bfseries Formula Sheet}\\[3pt]
  {\small\itshape Use this page only if your exam provides a formula sheet.}\\[2pt]
  \rule{\textwidth}{0.4pt}
\end{center}

\subsection*{Fitness}
\begin{center}\renewcommand{\arraystretch}{2.1}
\begin{tabular}{r@{\;}l@{\qquad}l}
  \(w\) & \(=\ \text{viability} \times \text{fecundity}\) & absolute fitness, offspring left per zygote\\
  \(W\) & \(=\ w / w_{\max}\) & relative fitness\\
  \(s\) & \(=\ 1 - W\) & selection coefficient
\end{tabular}\end{center}

\subsection*{Mean fitness and the general equation}
\begin{align*}
  \overline{W} &= f(AA)\,W_{AA} + f(Aa)\,W_{Aa} + f(aa)\,W_{aa}\\[4pt]
  q' &= \frac{\tfrac{1}{2}\,f(Aa)\,W_{Aa} + f(aa)\,W_{aa}}{\overline{W}}, \qquad \Delta q = q' - q
\end{align*}
Under Hardy--Weinberg, \(f(AA)=p^2\), \(f(Aa)=2pq\), \(f(aa)=q^2\). With counted individuals that selection acts on directly, use the counted frequencies.

\subsection*{Average excess}
\begin{align*}
  \alpha_a &= q\,W_{aa} + p\,W_{Aa} - \overline{W}, & \Delta q &= \frac{q}{\overline{W}}\,\alpha_a\\[4pt]
  \alpha_A &= p\,W_{AA} + q\,W_{Aa} - \overline{W}, & \Delta p &= \frac{p}{\overline{W}}\,\alpha_A
\end{align*}
Requires Hardy--Weinberg genotype frequencies.

\subsection*{Special models (\(q\) is always the allele selected against)}
\begin{align*}
  \text{Against the recessive:}\quad \Delta q &= \frac{-s\,p\,q^2}{1 - s\,q^2}\\[4pt]
  \text{Against the dominant:}\quad \Delta q &= \frac{-s\,p^2 q}{1 - 2sq + sq^2}
\end{align*}

\subsection*{Heterozygote advantage or disadvantage}
\[
  \hat{q} = \frac{s_1}{s_1 + s_2}
\]
with \(s_1\) for the \(p\) homozygote and \(s_2\) for the \(q\) homozygote. The equilibrium is stable under heterozygote advantage and unstable under heterozygote disadvantage.

\label{LastPage}
\end{document}